\documentclass[11pt]{article}
\usepackage[a4paper,margin=1in]{geometry}
\usepackage{fontspec}
\usepackage[boldfont=false]{xeCJK}
\setCJKmainfont{FandolSong-Regular.otf}
\usepackage{amsmath}
\usepackage{amssymb}
\usepackage{array}
\usepackage{booktabs}
\usepackage{graphicx}
\usepackage{adjustbox}
\usepackage{enumitem}
\setlist[itemize]{topsep=2pt,partopsep=0pt,parsep=0pt,itemsep=2pt,leftmargin=1.4em}
\setlist[enumerate]{topsep=2pt,partopsep=0pt,parsep=0pt,itemsep=2pt,leftmargin=1.6em}
\usepackage{parskip}
\usepackage{fvextra}
\fvset{breaklines=true,breakanywhere=true,fontsize=\small}
\setlength{\parindent}{0pt}

\begin{document}

\begin{center}
  {\LARGE\bfseries Motion in a circle}\\[0.35em]
  {\large A Level Physics}
\end{center}
\vspace{0.6em}

\section*{Angles in radians}

The \textbf{radian} 弧度 is the angle made at the centre of a circle by an \textbf{arc} 弧 whose length equals the radius. For an arc of length $s$ on a circle of radius $r$, the angle in radians is

\[
\theta = \frac{s}{r}.
\]

Radians have no unit (a ratio of lengths). A full circle has $s = 2\pi r$, so $\theta = 2\pi\ \text{rad}$. A half-circle is $\pi\ \text{rad}$; a quarter is $\pi/2\ \text{rad}$.

\par\medskip
\noindent\begin{minipage}{\linewidth}
\centering
\adjustbox{max width=\linewidth, max totalheight=0.42\textheight}{%
  \includegraphics{/Users/joschka/Desktop/Programming/eXercise/output/eXam/handouts/a_level_physics/assets/12-p239-radian-definition.pdf}}\\[0.3em]
\emph{One radian is the angle whose arc length equals the radius}
\end{minipage}\par\medskip

To convert: $1\ \text{rad} = 180°/\pi \approx 57.3°$. Set your calculator to \textbf{radians} for this topic; "degree" mode will give wrong answers.

\section*{Uniform circular motion: angular speed}

An object moves in a circle of radius $r$ at constant speed $v$. Define:

\begin{itemize}
\item \textbf{angular displacement} 角位移 $\theta$ — the angle (in radians) turned through by the radius from a chosen start line.

\item \textbf{angular speed} 角速度 $\omega$ — the rate of change of angular displacement.

\end{itemize}

For uniform motion $\omega$ is constant and

\[
\omega = \frac{\theta}{t}.
\]

Unit: $\text{rad s}^{-1}$.

\subsection*{Period and frequency}

If the object goes once round ($2\pi\ \text{rad}$, one \textbf{revolution} 圈) in time $T$ (the \textbf{period} 周期), then

\[
\omega = \frac{2\pi}{T} = 2\pi f,
\]

where $f = 1/T$ is the \textbf{frequency} 频率 of turning (Hz).

\subsection*{Linear and angular speed}

In one period $T$ the object travels a distance $2\pi r$ (the \textbf{circumference} 周长) at constant speed, so

\[
v = \frac{2\pi r}{T} = r \omega.
\]

This links the linear (\textbf{tangential} 切向) speed $v$ with the angular speed $\omega$. At a larger radius (for the same angular speed) the linear speed is larger — a child on the edge of a merry-go-round moves faster than one near the centre, even though both go round once in the same time.

\par\medskip
\noindent\begin{minipage}{\linewidth}
\centering
\adjustbox{max width=\linewidth, max totalheight=0.42\textheight}{%
  \includegraphics{/Users/joschka/Desktop/Programming/eXercise/output/eXam/handouts/a_level_physics/assets/12-p240-angular-speed.pdf}}\\[0.3em]
\emph{As the radius turns through Δθ the object moves an arc Δs at speed v}
\end{minipage}\par\medskip

\section*{Centripetal acceleration}

An object moving in a circle at constant speed still has a \textbf{changing velocity} 速度 — its direction keeps changing, even though its size stays the same. A changing velocity needs an \textbf{acceleration} 加速度. This acceleration points \textbf{towards the centre} and is the \textbf{centripetal acceleration} 向心加速度.

\subsection*{Size}

\[
a = \frac{v^{2}}{r} = r\omega^{2}.
\]

The two forms are equal because $v = r\omega$. Pick the one with the quantities you have.

The centripetal acceleration is \textbf{perpendicular} 垂直 to the velocity at every instant — never along the direction of motion. (If part of it were along the motion, the speed would change.) Unit: $\text{m s}^{-2}$.

\par\medskip
\noindent\begin{minipage}{\linewidth}
\centering
\adjustbox{max width=\linewidth, max totalheight=0.42\textheight}{%
  \includegraphics{/Users/joschka/Desktop/Programming/eXercise/output/eXam/handouts/a_level_physics/assets/12-p241-centripetal-direction.pdf}}\\[0.3em]
\emph{The velocity points along the tangent; the force and acceleration point to the centre}
\end{minipage}\par\medskip

\section*{Centripetal force}

By Newton's second law, the resultant \textbf{force} 力 on a body in circular motion at constant speed is

\[
F = m a = \frac{m v^{2}}{r} = m r \omega^{2}.
\]

This is the \textbf{centripetal force} 向心力. It \textbf{always points towards the centre} — perpendicular to the velocity.

The centripetal force is not a new kind of force — it is the \textbf{net result} of the real forces acting (tension, gravity, friction, electric attraction, normal contact force, …). In a problem, work out which real force(s) provide it.

\subsection*{Where the centripetal force comes from}

\begin{itemize}
\item \textbf{Ball on a string} in a horizontal circle: the \textbf{tension} 张力 in the string.

\item \textbf{Car turning a flat corner}: the \textbf{friction} 摩擦力 between tyres and road ($F = m v^{2}/r$). If the car goes too fast, friction is not enough and it skids outwards.

\item \textbf{Banked corner} 倾斜 (no friction): the horizontal part of the \textbf{normal contact force} 支持力; $\tan\theta = v^{2}/(rg)$ for the angle that needs no friction.

\item \textbf{Planet or satellite} 卫星 in \textbf{orbit} 轨道: the \textbf{gravitational attraction} 引力, $G M m / r^{2} = m v^{2}/r$.

\item \textbf{Electron} 电子 in a circular orbit (Bohr-style model): the \textbf{electrostatic} 静电 attraction between the electron and the positive \textbf{nucleus} 原子核:

\end{itemize}

\[
\frac{kZe^{2}}{r^{2}} = \frac{m_{e} v^{2}}{r},
\]

where $k = 1/(4\pi\varepsilon_{0})$ and $Z$ is the nuclear charge. Solve for $v$ to get the orbital speed; then $T = 2\pi r/v$.

\par\medskip
\noindent\begin{minipage}{\linewidth}
\centering
\adjustbox{max width=\linewidth, max totalheight=0.42\textheight}{%
  \includegraphics{/Users/joschka/Desktop/Programming/eXercise/output/eXam/handouts/a_level_physics/assets/12-p243-banked-corner.pdf}}\\[0.3em]
\emph{On a banked track the horizontal part of the road's force provides the centripetal force}
\end{minipage}\par\medskip

\subsection*{Vertical circles}

When the circle is upright, the speed is \textbf{not} constant (gravity does work) — but at each instant the \textbf{net force towards the centre} still equals $m v^{2}/r$:

\begin{itemize}
\item at the bottom of a loop: tension up, \textbf{weight} 重力 down, so $T - mg = m v^{2}/r$ — the tension is largest here.

\item at the top of a loop: tension and weight both point down (towards the centre), so $T + mg = m v^{2}/r$ — the tension is smallest. For the slowest speed at the top with the string just tight, set $T = 0$: $mg = m v_{\text{min}}^{2}/r$, giving $v_{\text{min}} = \sqrt{gr}$.

\end{itemize}

\par\medskip
\noindent\begin{minipage}{\linewidth}
\centering
\adjustbox{max width=\linewidth, max totalheight=0.42\textheight}{%
  \includegraphics{/Users/joschka/Desktop/Programming/eXercise/output/eXam/handouts/a_level_physics/assets/12-p244-vertical-circle-forces.pdf}}\\[0.3em]
\emph{Forces on a person at the top and bottom of a vertical circle}
\end{minipage}\par\medskip

The constant-speed result ($v = r\omega$, $\omega$ constant) holds for horizontal circles, or where the force only bends the path (orbits in gravity, charges in a \textbf{magnetic field} 磁场).

\section*{How to structure a circular-motion answer}

\begin{enumerate}
\item \textbf{Find the radius} $r$ and choose $v$ or $\omega$. Use $v = r\omega$ to switch between them.

\item \textbf{Find the centripetal acceleration} with $a = v^{2}/r$ or $r\omega^{2}$.

\item \textbf{List the real forces} and write Newton's second law in the \textbf{radial} 径向 direction (towards the centre is positive). Set the net inward force equal to $m v^{2}/r$.

\item For \textbf{period or frequency}: use $\omega = 2\pi/T$, or $T = 2\pi r / v$.

\item Check the \textbf{directions}: centripetal force and acceleration point to the centre; the velocity is along the tangent.

\end{enumerate}


\end{document}
